\begin{solapartat}
La segona llei de Newton estableix que: $F = ma$ \punts{0,2}

Segons la llei de gravitació universal, el mòdul de la força és: $F = G\,\dfrac{mM}{r^2}$ \punts{0,2}

Per tant obtenim que: $a = G\,\dfrac{M}{r^2}$ \punts{0,1}

D'altra banda, considerant que els planetes descriuen moviments circulars uniformes, la seva velocitat es pot expressar com el perímetre de l'òrbita dividit pel període: $v = \dfrac{2\pi r}{T}$ \punts{0,1} i tenint en compte que en aquest cas l'acceleració centrípeta val: $a = \dfrac{v^2}{r}$ \punts{0,1}\unskip, obtenim que: $a = \dfrac{4\pi^2 r}{T^2}$ \punts{0,1}\unskip.

Per tant: $G\,\dfrac{M}{r^2} = \dfrac{4\pi^2 r}{T^2}$ que reordenem per obtenir la tercera llei de Kepler:
\[ \frac{r^3}{T^2} = \frac{GM}{4\pi^2} \Rightarrow T^2 \propto r^3 \qquad \text{\punts{0,2}} \]
\end{solapartat}

\begin{solapartat}
Aïllem $M$:

\punts{0,6} $M = \dfrac{4\pi^2 r^3}{GT^2} = \dfrac{4\pi^2\left(1,496\times 10^{11}\right)^3}{6,67\times 10^{-11}\times\left(3,154\times 10^{7}\right)^2} = 1,99\times 10^{30}\ \mathrm{kg}$ \punts{0,4}
\end{solapartat}
